\section{The bounded quadratic-seven family on the original window}
\label{sec:quadratic-seven-coefficient-two-window}

\paragraph{The complete preceding union and the same host.}
Keep every datum, coefficient completion, original lattice, page,
prescribed pair and local condition of AC-G188-window.
In particular,
\[
 \begin{gathered}
 H_u=\mathbb Q(\zeta_{2^{h(u)+2}},\zeta_7,
           \sqrt{-q},\sqrt{-d},\omega,\sqrt[3]{269}),
             \qquad K=\mathbb Q(\omega),\\
 u^{2+8\varepsilon}\leq2^{h(u)}\leq2u^{2+8\varepsilon},
 \qquad0<\varepsilon<1/256,
 \qquad T=\lfloor u^{399}\rfloor.
 \end{gathered}\tag{Q7W.1}
\]
Write $Z^{\rm three,-8}$ and $I_{399}^{\rm three,-8}$ for the
exact preceding omitted set and remainder in (N3W.2)--(N3W.5).
Retain its notation
$N_T=T-1$, $C(h),E_3(h),P(h),B_T,A_T,R_T,S_T,J_-(T)$,
including the unchanged original residue embedding defining $E_3(h)$.
Let $\xi=\chi_{-7}$ be the unique quadratic character modulo seven.

For integers $h\geq2$ and $T\geq1$, put
\[
 \begin{aligned}
 Z^{\rm quad7}_2(u)
   &=\{(2,(\delta\xi)|_{G_K},e):
       \delta\in\widehat{(\mathbb Z/2^{h+2}\mathbb Z)^\times},
                          \ 1\leq e<T\},\\
 Z^{\rm all7}_2(u)&=Z^{\rm three,-8}(u)\sqcup Z^{\rm quad7}_2(u).
 \end{aligned}\tag{Q7W.2}
\]
All auxiliary quadratic and Kummer cubic factors in the new labels
are trivial.  Define $I_{399}^{\rm all7,2}$ by omitting exactly
this union from the original positive-index sum (TC.1).
The marker two records the new family; the union retains all
previous omissions at every coefficient prime.

\begin{theorem}[AC-G190-window: the complete seven-factor family at two]
\label{AC-G190-window}
Every new triple in (Q7W.2) is an actual character and weight
on its unchanged original coefficient $O_2$.
The displayed union is disjoint, with
\[
 \begin{aligned}
 |Z^{\rm all7}_2\setminus Z^{\rm three,-8}|
       &=2^{h+1}N_T,\\
 |(Z^{\rm all7}_2)_{\ell=2}|&=6\cdot2^{h+1}N_T,\\
 |Z^{\rm all7}_2|
       &=(10+C(h)+E_3(h)+P(h)+6\cdot2^{h+1})N_T\\
       &\hspace{1cm}-B_T+12A_T-R_T-S_T+J_-(T).
 \end{aligned}\tag{Q7W.3}
\]
Thus the coefficient-two part contains every product of a
two-power character with every character modulo seven, at every
positive weight below $T$, with the other factors trivial.
All counts at the other coefficient primes are unchanged.
The omitted set is properly enlarged for $T\geq2$.
No strict decrease of the numerical remainder is required.

Let $Q_u$ be the normalized positive derived-index cost of the
\emph{entire} union (Q7W.2), and set
\[
 \mathcal E_u(T)=8D_u^2T^2\log(8TD_u)
        +\tfrac12D_u\mathcal R_uT^2+2D_ur_uc_0T.
 \tag{Q7W.4}
\]
There is an absolute constant $C$ such that
\[
 \begin{aligned}
 I_{399}&=I_{399}^{\rm all7,2}+Q_u,\\
 I_{399}^{\rm all7,2}&\leq I_{399}^{\rm three,-8}\leq I_{399},\\
 0\leq Q_u&\leq\mathcal E_u(T)
              +C2^hT^2\bigl(h+\log(T+2)+1\bigr)\\
 &=O_\varepsilon\bigl(u^{802+16\varepsilon}(\log u)^2\bigr).
 \end{aligned}\tag{Q7W.5}
\]
The first term is one complete (FL.4) error for the full union.
The second is the actual strict-group cost allowed by AC-G189;
its vanishing is not assumed.  It is not another local-error bound.
These statements concern nonnegative extended sums; no individual
strict group or positive derived index in the new family is
asserted zero.

With the exact (TC.3) quantifiers and positive relative slack,
a numerical bound for $I_{399}^{\rm all7,2}$ implies the original
full numerical premise on the identical common window.
Every original source parameter, including the actual rank $R.t$,
is retained.
\end{theorem}

\begin{proof}
\proofstep{1. Literal application of the original-integral bound.}
The actual host contains $\mathbb Q(\zeta_{2^{h+2}},\zeta_7)$.
Every $\delta\xi$ in (Q7W.2) is therefore an actual finite label.
The exponent of the two-power character group is $2^h$ for
$h\geq2$.  The quadratic factor has order two, so all values of
$\delta\xi$ already belong to the prescribed completion $O_2$
of $\mathbb Z[\zeta_{3\cdot2^h}]$.
With $E_2=\operatorname{Frac}(O_2)$ the coefficient in (FC.2) is
canonically
\[
 (O_2\otimes_{\mathbb Z_2}\mathbb Q_2/\mathbb Z_2)
                     (\delta\xi)(-e)
             =(E_2/O_2)(\delta\xi c_2^{-e}).
 \tag{Q7W.6}
\]
Apply (G189.1)--(G189.4), with this original ring and the same
positive integer $e$.  No extra continuous character is inserted.
The theorem allows all original ramification and divides the
logarithmic cardinality by $b_2=[E_2:\mathbb Q_2]$, exactly
the normalization in (FC.2).

The strict conditions are full local zero above two and zero
in every other finite inertia group.  The full strict comparison
in (G189.19) and its proof injects the original (FC.2) group
into that absolute strict group by inflation.
Outside the original ramification set, an inflated cocycle kills
inertia already.  At every original or added host prime the
specified condition remains in place.  Thus every actual group
in (Q7W.2) is bounded, with
\[
 \sum_{(2,\delta\xi,e)\in Z^{\rm quad7}_2}
          \frac{\log\#\mathcal S_{2,\delta\xi}(e)}{b_2}
       \leq C2^hT^2\bigl(h+\log(T+2)+1\bigr).
 \tag{Q7W.7}
\]
Every old triple retains its original strict zero from G188.
The new inequality, rather than a new vanishing assertion, is
what is used below.  In particular it is consistent with the
nonzero strict classes of G186.

\proofstep{2. Disjointness, exact counts and the complete seven factors.}
The preceding coefficient-two set consists exactly of
$\delta\epsilon$ at every weight, where $\epsilon$ modulo seven
is trivial or has order three or six.
There are five such characters: one trivial, two of order three,
and two of order six.  The sole missing character is $\xi$ of
order two.  No other coefficient-two family was added in G188.

All products $\delta\epsilon$ in this discussion are distinct
over $\mathbb Q$ by the coprime conductors two and seven.
They remain distinct over $K$: a rational character trivial on
$G_K$ is either $1$ or $\chi_{-3}$, and a ratio ramified only
at two and seven cannot equal the conductor-three character.
The other old coefficient primes distinguish their triples
regardless of any finite-character identity.
This proves the disjointness in (Q7W.2).

There are $\varphi(2^{h+2})=2^{h+1}$ two-power characters and
$N_T=T-1$ positive integer weights below $T$.
This gives the first count in (Q7W.3); adjoining the missing
sixth seven-factor gives its second count.
Adding $2^{h+1}N_T$ to the exact cumulative count (N3W.4)
gives the third count, without altering any old residue exception
or count at another coefficient prime.
When $T\geq2$, the new triple with $\delta=1$ and $e=1$ is
present, proving proper inclusion of omitted sets.
Its inclusion is not justified by claiming its strict group zero.

\proofstep{3. One full-union local error and the bounded strict cost.}
Keep the original codifferent duality in (FL.2) and the original
$b_\ell=\operatorname{rank}_{\mathbb Z_\ell}O_\ell$.
For each triple in the entire union, the determinant comparison
(EU.2) and its all-prime application give
\[
 \begin{aligned}
 \frac{\log\#k_{O_\ell}}{b_\ell}\max(0,i_{\ell,\eta,e})
 &\leq \frac{\log\#H^2(C_{\ell,\eta,e})}{b_\ell}\\
 &\leq\frac{\log\#\mathcal S_{\ell,\eta}(e)}{b_\ell}
   +\frac{\bigl|\log\#H^2(C_{\ell,\eta,e})
                 -\log\#\mathcal S_{\ell,\eta}(e)\bigr|}{b_\ell}.
 \end{aligned}\tag{Q7W.8}
\]
The first summand is zero on the old union and is bounded by
(Q7W.7) in total on the new family.
At each weight the second summand, over the \emph{entire} union,
is one nonnegative subsum of the full (FL.4) sum over all
coefficient primes and all actual characters.
Sum that full error once for $1\leq e<T$, using
$\sum_{1\leq e<T}e\leq T^2/2$, to obtain (Q7W.4).
This proves the stated upper bound for $Q_u$.
There is no multiplication of the full error by the number of
characters or by the number of old subfamilies.
The original integral Fourier corrections, in particular at
two and three, remain in the unchanged source interface.

The existing bounds $D_u\leq288u^{2+8\varepsilon}$,
$r_u\leq12$ and
$\mathcal R_u=O_\varepsilon(u^{\varepsilon/5}(\log u)^2)$
give the old polynomial envelope for $\mathcal E_u(T)$.
Also $2^h\leq2u^{2+8\varepsilon}$ and
$h+\log(T+2)+1=O_\varepsilon(\log u)$ on sufficiently large
source parameters, so (Q7W.7) is
$O_\varepsilon(u^{800+8\varepsilon}\log u)$.
It is covered by the larger envelope in (Q7W.5), uniformly
in every prescribed pair and original coefficient.
Equivalently it is bounded by an absolute multiple of
$D_u^2T^2\log(8TD_u)$, since $D_u=144\cdot2^h$.
The additive equality in (Q7W.5) follows by partitioning
nonnegative sums; it does not subtract two infinities.

\proofstep{4. The identical-window quantifiers.}
Fix any retained finite datum, page and $0<\varepsilon<1/256$.
Assume exactly the (TC.3) premise for $I_{399}^{\rm all7,2}$:
there exist $J\geq1$, $X_0$ and $s\in(0,1)$, independent of
$X,q,d,u,\ell,\eta,e$, such that every prescribed pair at
$X\geq X_0$ has one $W$, with $X^{1/J}\leq W$ and $W^2\leq X$,
on which the new remainder is at most
$(1-s)\exp(u^\varepsilon/2)$ for every $u\in[W,W^2]$.
The uniform error in (Q7W.5) is
$o(\exp(u^\varepsilon/2))$.
Choose $U_0\geq2$ making it at most
$(s/2)\exp(u^\varepsilon/2)$ for every permitted $u\geq U_0$,
and increase the threshold to $\max(X_0,U_0^J)$.
Keep exactly the supplied $W$.
It satisfies $W\geq U_0$, so the original full numerical premise
holds throughout that same interval with slack $s/2$.

Every preceding remainder premise implies the new one by
nonnegativity.  Uniform little-$o$ premises and fixed finite
simultaneous families use the same proof with a maximum of
their uniform constants; no separate windows are intersected.
The original field, coefficient orders, local conditions,
prescribed pair and actual $R.t$ remain unchanged through K35.
The retained mixed theta/Selmer realization
$\mathcal R_{\rm cub}$ is still an input with its exact original
coefficient and local normalizations.  No AW.1 premise is added.
\end{proof}

\paragraph{Boundary and retained dictionary.}
Only the previously missing pure quadratic-seven coefficient-two
family has been added to the omitted set, at its proved bounded
cost.  This gives the complete six-seven-factor union at two;
it does not cover nontrivial cubic or auxiliary quadratic factors.
Every other unsupported combination remains in the remainder.
The all-prime numerical bound, mixed realization and H04.4
remain open.  Neither G191's critical-component bound nor
G193's noncritical lemma is used in this composition.
The previously audited SGA~1, Expos\'e~V, \S\S4,7 and Stacks~03PK
finite-cover/Artin--Tate dictionary is retained; no new archive
attribution, mixed motive or Lean execution is claimed.
